\documentclass[11pt]{article}
\usepackage[T1]{fontenc}
\usepackage[utf8]{inputenc}
\usepackage[french]{babel}
\usepackage{newpxtext,newpxmath}
\usepackage[a4paper,top=1.6cm,bottom=1.55cm,left=2.1cm,right=2.1cm]{geometry}
\usepackage{microtype}
\usepackage{xcolor}
\usepackage{titlesec}
\usepackage[most]{tcolorbox}
\usepackage{fancyhdr}
\usepackage{enumitem}
\usepackage{booktabs}
\usepackage{listings}
\lstdefinestyle{jsonic}{basicstyle=\ttfamily\footnotesize,backgroundcolor=\color{boxgray},%
  breaklines=true,columns=fullflexible,keepspaces=true,xleftmargin=6pt,%
  framexleftmargin=6pt,aboveskip=4pt,belowskip=4pt}
\lstset{style=jsonic}

\definecolor{efenavy}{HTML}{1F3A5F}
\definecolor{efeteal}{HTML}{0E7A6E}
\definecolor{efeamber}{HTML}{B25A1E}
\definecolor{efeplum}{HTML}{7A4EA8}
\definecolor{efegray}{HTML}{54606E}
\definecolor{boxgray}{HTML}{F3F5F8}
\definecolor{boxamber}{HTML}{FBF3EA}
\definecolor{boxedge}{HTML}{D9E0EA}

\titleformat{\section}{\color{efenavy}\large\bfseries}{}{0pt}{}
\titlespacing{\section}{0pt}{7pt}{3pt}
\titleformat{\subsection}[runin]{\color{efenavy}\bfseries}{}{0pt}{}[\;---\;]
\titlespacing{\subsection}{0pt}{5pt}{0.5em}

\setlength{\parskip}{0.6pt}
\setlength{\parindent}{0pt}
\setlength{\emergencystretch}{3em}

\pagestyle{fancy}\fancyhf{}
\renewcommand{\headrulewidth}{0.3pt}
\lhead{\small\color{efegray}Scella}
\rhead{}
\cfoot{\small\color{efegray}\thepage}

\newcommand{\mesure}[1]{{\color{efeteal}\textbf{Mesuré.} }#1\par}
\newcommand{\concept}[1]{{\color{efenavy}\textbf{Concept —} }#1\par}
\newcommand{\exemple}[1]{{\color{efeamber}\textbf{Exemple —} }#1\par}
\newcommand{\utilite}[1]{{\color{efeplum}\textbf{Utilité réelle —} }#1\par}

\newtcolorbox{these}{colback=boxgray,colframe=boxedge,boxrule=0.5pt,arc=2.5pt,
  left=8pt,right=8pt,top=6pt,bottom=6pt}
\newtcolorbox{cadre}{colback=boxamber,colframe=efeamber!55,boxrule=0.6pt,arc=2.5pt,
  left=8pt,right=8pt,top=6pt,bottom=6pt}

\begin{document}

\begin{center}
{\Huge\bfseries\color{efenavy} Scella}
\end{center}
\vspace{6pt}

\begin{these}
\textbf{Le problème.} Faire coopérer plusieurs IA (Claude, GPT, un modèle interne) sur une même
décision en milieu réglementé, et prouver après coup — à un tiers, sans lui confier ses données ni
lui demander de croire l'opérateur sur parole — \emph{exactement} ce que chacune a fait.

\textbf{Ce que Scella résout concrètement :}
\begin{itemize}[leftmargin=1.3em,itemsep=1.5pt,topsep=2pt,parsep=0pt]
\item orchestrer \textbf{plusieurs fournisseurs d'IA} sans glue tierce, avec une trace auditable ;
\item garder les \textbf{données chez leur propriétaire} (multi-région) tout en coordonnant l'IA ;
\item rendre chaque décision \textbf{rejouable et re-vérifiable hors-ligne}, sans accès aux systèmes ;
\item \textbf{réduire coût et hallucination} sur les documents longs : compresser avant le gros
modèle, vérifier après contre les sources.
\end{itemize}
\vspace{2pt}
\textbf{L'interstice prend trois formes :} \textbf{L'Encadrement} — compresser le prompt avant
l'appel, vérifier la réponse après (le simple au SLM, le raisonnement au gros modèle) ;
\textbf{l'essaim médiateur SLM} — l'arête entre deux LLM devient une chaîne déclarée de petits
agents scellés (traduire → compresser → vérifier) ; et, \textbf{en parallèle d'un agent de code},
de petits agents traitent le flux d'un gros modèle à côté (vérifier contre le dépôt, contrôler un
schéma) sans bloquer sa boucle. Six briques, deux lignées, détaillées ci-dessous.
\end{these}

\section{I \;\textbar\; La lignée latente : nano-robots, essaim, méta-essaim}

\subsection{Nano-robots}
Le nano-agent est l'\textbf{unité élémentaire} de la lignée latente : une \emph{ligne d'un tenseur}, un
vecteur dans l'espace de \texttt{efe-core}, mis à jour par un pas de couplage. Sa dynamique est
\textbf{déterministe} et ses invariants sont explicites dans le code : population fixe, priors
$\sin(i)$ neutres, une seule ronde de fusion. C'est l'unité que l'essaim agrège et que le méta-essaim distribue.
\concept{Un nano-agent = une ligne du tenseur d'état, mise à jour par un pas de couplage fixe.
Pas un processus : un vecteur de nombres.}
\utilite{L'unité adressable qui rend chaque calcul \emph{traçable et rejouable à l'identique} :
base de l'audit et de la reproductibilité exigée en environnement réglementé.}
\exemple{Sur une question, on tire 16 échantillons du modèle : chacun est une ligne du tenseur
d'état, l'agent n$^\circ$3 = le 3\ieme échantillon. Le pas de couplage les rapproche.
\emph{Autre cas :} 16 échantillons déjà d'accord $\rightarrow$ un pas ne change rien ; dispersés
$\rightarrow$ il les resserre.}

\subsection{L'essaim ($\Phi$)}
L'essaim \textbf{fusionne les nano-agents par un point-fixe latent} $\Phi$ : une ronde,
déterministe (\texttt{init\_sigma}=0, graine dérivée de la spec), qui condense la population en
un \textbf{centroïde de consensus} : la position latente vers laquelle la population converge en
une ronde, reproductible bit à bit et donc \textbf{rejouable} — on re-déroule une fusion et
l'on retrouve le même résultat, ce qui alimente l'audit. Le déterminisme est durci jusqu'à rendre
$\Phi$ \emph{pur par construction} : deux fusions concurrentes n'entrent jamais en conflit, une
même spec produit toujours le même résultat.
\mesure{Le point-fixe est caractérisé par un \textbf{estimateur de Lipschitz adverse} (par
bissection, sans accès interne au modèle) qui en donne la constante empirique —
$\hat{L}\approx0{,}69$ en topologie dynamique, $0{,}37$ en topologie fixe — et sert de
\textbf{garde fail-closed}. Une passe de $\Phi$ coûte $\sim$2,6 s (zérotage d'une arène de
1 GiB).}
\concept{Fusionner = itérer le couplage jusqu'à un point-fixe ; la sortie est le centroïde vers
lequel les lignes convergent.}
\utilite{Un résultat déterministe assorti d'une \emph{borne de stabilité mesurée} (Lipschitz) —
le type de borne qu'une certification de sûreté peut exiger (automobile, médical, aéronautique).}
\exemple{Les 16 échantillons d'une réponse convergent en une ronde vers leur position
consensuelle (le centroïde), reproductible bit à bit. \emph{Autre cas :} 12 échantillons
proches et 4 aberrants $\rightarrow$ le centroïde est tiré vers les aberrants (moyenne uniforme,
pas de détection d'aberrant).}

\subsection{Le méta-essaim}
Le méta-essaim \textbf{coordonne plusieurs essaims} par gRPC. C'est la \emph{couche de
coordination} de la lignée latente, réelle et testée : participants purs
Tensor$\rightarrow$Tensor, \textbf{manifeste de moteur transmis dans le protocole}
(\texttt{kappa} et graine explicites, si bien qu'un binaire d'une
autre version refuse une ligne moteur qu'il ne connaît pas), \textbf{un contrôle strict au
démarrage}, et \textbf{mTLS strict} (TLS 1.3, vérification SAN). La pureté des participants et le sceau de bout en bout en
font la \textbf{base de déploiement (C.3)} : l'infrastructure qui \emph{exécutera} les hops de
la lignée interstice, sous les mêmes garanties d'intégrité que le reste du système. Le graphe
d'exécution se \textbf{re-dérive depuis la topologie déclarée}, ce qui rend cette couche
\emph{auditable} indépendamment de ce que les participants calculent.
\concept{Plusieurs essaims sur des machines séparées ; un coordinateur échange leurs centroïdes
en gRPC + mTLS et scelle chaque envoi.}
\utilite{Faire tourner l'IA \emph{là où sont les données} (multi-région, multi-tenant) avec une
coordination scellée et auditable : résidence des données, souveraineté (banque, santé, public).}
\exemple{Plusieurs machines traitent chacune leur lot en parallèle ; le coordinateur échange
leurs résultats sous mTLS et scelle chaque envoi ; un audit re-trace qui a produit quoi, sans
rejouer les calculs. \emph{Autre cas :} une machine tourne une version incompatible
$\rightarrow$ sa ligne moteur inconnue la fait refuser.}

\section{II \;\textbar\; La lignée interstice : Le Relais inter-modèles, L'Encadrement}

\subsection{Le Relais inter-modèles}
Le Relais inter-modèles intercale \textbf{un petit modèle} (Qwen3-0.6B, décodage contraint par grammaire GBNF)
entre deux appels frontière : il traduit le dialecte de l'un (\texttt{ClaudeXml}) vers une
forme canonique, puis vers le dialecte de l'autre (\texttt{Gpt4Json}). C'est la brique la
plus \textbf{aboutie} : codée et poussée de bout en bout, servie en
opt-in, protégée par un disjoncteur à sonde unique, avec une \textbf{provenance par-hop
scellée et re-vérifiable} (chaque hop est un n\oe ud HMAC ; \texttt{slm\_mediated} n'est vrai
que si \emph{tous} les hops sont SLM) et des \textbf{deadlines de mission qui bornent le hop}
(dégradation estampillée, six raisons stables). Construit en cinq jalons (seam byte-identique
quand le SLM est absent, grammaire, service opt-in, provenance chaînée, deadlines), avec un
audit de complétude sans trou et quatre revues adverses (vingt-six correctifs).
Le Relais se justifie par sa \textbf{fonction} : traduction de dialectes, provenance scellée,
souveraineté — chaque échange inter-modèle devient un hop \emph{attribuable et re-vérifiable},
relayé d'un modèle à l'autre. C'est la brique dont toute la lignée interstice hérite ce point d'intégration.
Le vocabulaire du petit modèle est fixé par une \textbf{grammaire d'union GBNF} qui lie chaque
schéma à sa forme, et six raisons de dégradation stables rendent chaque repli lisible dans le
reçu.
\concept{Entre deux gros modèles, un petit modèle réécrit la sortie de l'un dans la forme exacte
qu'attend l'autre, et scelle cette réécriture.}
\utilite{Faire coopérer plusieurs fournisseurs d'IA (Claude, GPT, modèle interne) \emph{sans
glue tierce}, avec une trace auditable : orchestration multi-fournisseur en entreprise réglementée
(finance, assurance, défense).}
\exemple{Claude extrait une clause d'un contrat fournisseur et la rend en XML ; l'étape
suivante attend du JSON pour la confronter à une politique. Le petit modèle traduit via une
forme canonique ; le reçu garde les hash. \emph{Autre cas :} Gemini rend la même clause en
Markdown $\rightarrow$ canonique $\rightarrow$ la forme XML qu'attend Claude.}

\subsection{L'Encadrement}
L'Encadrement généralise le Relais : au lieu de seulement \emph{traduire entre} deux modèles, il
\emph{encadre} un appel frontière — un \textbf{compresseur} avant, un \textbf{vérificateur}
après — dans le même reçu re-vérifiable. \textbf{Le principe : au petit modèle des tâches
\emph{simples}} — compresser, confronter à des sources — \textbf{au gros modèle le
raisonnement.} Les petits agents sont bornés, contraints par grammaire, et \textbf{agnostiques
du LLM} qu'ils encadrent : le même encadrement sert autour de \textbf{Claude, GPT ou Gemini}.
C'est additif : un type de hop
\{Prep, Translate, Check\} où un hop hérité sans étiquette reste un Translate, sceaux intacts.
Le compresseur (Llama-3.2-1B, température 0, \emph{map-reduce au-delà de $\sim$512 tokens})
réduit le prompt ; le vérificateur (Phi-4-mini, verdict fermé) annote une assertion — un
\textbf{ajout au reçu}, en parité stricte avec le chemin non-vérifié. Le \texttt{debate} est
\emph{exclu} du prep (dérive sémantique du sujet) ; les seuils sont \emph{réutilisés}, non
recopiés, depuis \texttt{efe-coordinator}. \texttt{slm\_mediated} reste gelé aux seuls hops de
traduction.
\textbf{Statut : architecture écrite ($\sim$535 lignes, jalons a$\rightarrow$d), prête à coder.}
Le compresseur et le vérificateur s'attachent \textbf{par rôle} (\texttt{prep\_for} /
\texttt{check\_for}), comme le traducteur du Relais, et écrivent dans le même reçu : l'Encadrement
est une extension naturelle du hop.
\concept{Autour d'un appel frontière : un petit modèle avant (compresse le prompt), un autre
après (vérifie la réponse) — les deux dans le même reçu.}
\utilite{Analyse de documents longs (juridique, assurance, achat) : réduire ce qu'on envoie au
modèle et vérifier sa réponse contre les sources — éviter une hallucination coûteuse sur un contrat.}
\exemple{Question sur un contrat de 40 pages : le compresseur réduit le contrat aux passages
pertinents avant l'appel ; le vérificateur confronte la réponse (« clause 7 : notification sous
72 h ») aux passages cités $\rightarrow$ SUPPORTED, en annotation. \emph{Autre cas :} la réponse
contredit la clause $\rightarrow$ CONTRADICTED, signalé, l'appel passe quand même.}

\section{III \;\textbar\; L'essaim médiateur SLM}

C'est un \textbf{autre type d'essaim} que celui de la Section~I : l'essaim latent
(nano $\rightarrow$ $\Phi$ $\rightarrow$ méta) est Tensor$\rightarrow$Tensor et
\emph{sélectionne} un consensus ; l'essaim médiateur — celui-ci — est fait de petits modèles
de langage qui \emph{transportent} l'information entre les gros.

\textbf{La thèse.} Une arête entre deux agents LLM est un interstice qu'on peut instrumenter.
Le Relais y avait mis un seul traducteur ; l'essaim médiateur y met une \textbf{chaîne déclarée de
petits agents spécialisés} — traduire, compresser, vérifier — chacun restant un \emph{hop
scellé}, borné en temps, attribuable. On ne produit pas « plus d'intelligence » : on produit un
reçu re-vérifiable de tout ce qui a transité.

\concept{Élargir le Relais : au lieu d'un petit modèle par arête, une chaîne déclarée (traduire
$\rightarrow$ compresser $\rightarrow$ vérifier $\rightarrow$ traduire), chaque étape un hop
scellé.}
\utilite{Pipelines de décision réglementés (achat, conformité, KYC, sinistres) : chaque étape
scellée, la décision arrive avec un reçu re-vérifiable — l'audit est intégré, pas ajouté après.}
\exemple{Demande d'achat : Claude rédige l'évaluation de risque du fournisseur $\rightarrow$
traduire $\rightarrow$ vérifier le fournisseur contre une liste de sanctions $\rightarrow$
compresser $\rightarrow$ GPT recommande une décision ; le reçu trace les 4 hops. \emph{Autre
cas :} le vérificateur de sanctions tombe $\rightarrow$ son hop dégrade, estampillé, la chaîne
continue et le reçu le dit.}

\textbf{En parallèle d'un LLM (agent de code).} Autre topologie : l'essaim tourne \emph{à côté} d'un
seul modèle, pas entre deux. Pendant qu'un agent écrit du code, de petits agents traitent son flux
en parallèle (vérifier une affirmation contre le dépôt, contrôler qu'une édition respecte un schéma)
sans bloquer sa boucle, chacun un hop scellé dans le reçu. \emph{Non conçue en détail, non codée.}

\subsection{Ce que c'est, précisément}
Une arête $A\rightarrow B$ n'est plus un hop unique mais une \emph{liste ordonnée} (traduire →
compresser → vérifier → retraduire), chaque étape portée par un petit modèle choisi \emph{par rôle}
dans un registre. La composition est \textbf{déclarée}, gelée par type de mission.

\subsection{Ce qui existe déjà, ce qui est neuf}
Environ \textbf{les deux tiers} sont là, hérités du Relais (hop, gardes, disjoncteur, sceau
par-hop, checker). Le neuf : un \emph{registre par rôle}, la chaîne \emph{typée}, et le
\emph{runtime} qui scelle $N$ n\oe uds.

\subsection{Les étapes d'un hop, et le reçu}
Chaque petit agent passe les gardes, appelle son modèle sous grammaire GBNF, valide la forme typée
(\emph{rejette} si fausse), s'estampille (moteur + raison), et se scelle en HMAC. La mission produit
\textbf{un seul reçu re-vérifiable} ; \texttt{slm\_mediated} n'est vrai que si \emph{tous} les hops
de traduction sont SLM, et falsifier une étiquette scellée casse la re-preuve.

\subsection{Le Flexible MLP — un filtre d'anomalie en amont (futur)}
Le \textbf{Flexible MLP} est un petit réseau quantifié (int8, quelques kilo-octets) qui
\emph{s'auto-adapte sur l'appareil} sous un \textbf{invariant d'adaptation prouvé} — il change sans
jamais violer sa garantie. Placé \emph{en amont} de la chaîne, il calcule un \textbf{score
d'anomalie} qui \emph{route} : les cas de confiance prennent la voie allégée, les cas anormaux vont
à la vérification complète (ou à l'humain). On ne paie le contrôle lourd que sur la fraction
anormale. \textbf{Sa limite} : il réduit le coût, il n'établit pas la confiance. \textbf{Pièce
séparée, pas encore intégrée à Scella.}

\vspace{3pt}
{\color{boxedge}\hrule height 0.6pt}
\vspace{3pt}
{\color{efenavy}\textbf{Le reçu — ce qui relie tout Scella.}} Les six briques n'écrivent qu'\textbf{un
seul reçu re-vérifiable} : chaque étape y est le même \textbf{n\oe ud HMAC scellé et chaîné},
re-dérivable par tout détenteur de la clé ; le \textbf{Flexible MLP}, une fois intégré, y scellera sa
décision de triage comme un n\oe ud de plus. \textbf{Garantie exacte —} \emph{tamper-evident pour qui
détient la clé et épingle le dernier maillon hors-bande} : pas \emph{trustless}, ancrage externe pas
encore branché, Flexible MLP pas encore intégré. Un actif de \textbf{discipline vérifiable}, pas un
coffre-fort distribué.

\clearpage
\section{Résultats -- ce que ça donne, par partie}
Ce que les mesures réelles ont établi de positif, brique par brique :
\begin{itemize}[leftmargin=1.3em,itemsep=2.5pt,topsep=3pt,label=\textendash]
\item \textbf{Nano-robots \& le reçu} -- socle \textbf{testé et relu} : près de \textbf{1\,200 tests
verts} sur le runtime, la chaîne d'audit et la provenance ; chaque étape scellée et re-vérifiable.
\item \textbf{Essaim $\Phi$ (fusion latente)} -- un petit modèle (3,8\,B) \emph{seul} \textbf{égale un
gros} (12\,B) sur une tâche réelle (\textbf{87,5\,\%}), à moindre coût ; le raisonnement le fait
passer de 40\,\% à \textbf{94\,\%} sur l'arithmétique (la bonne réponse \textbf{présente dans les
candidats à 100\,\%}) ; exécution \textbf{déterministe, rejouable bit à bit}, sous garde de
stabilité mesurée.
\item \textbf{Méta-essaim (coordination)} -- validé \emph{en live} sur de vrais fournisseurs (Claude
+ Gemini) : \textbf{coût réduit} ($\sim$\,−27\,\% agrégé, jusqu'à $\sim$90\,\% sur le routable),
\textbf{qualité tenue}, latence jusqu'à divisée par $\sim$3 sur les cas rapides.
\item \textbf{Le Relais (Fil B)} -- \textbf{servi et complet} (B.0→B.4), durci par revue adverse
($\sim$19 correctifs), $\sim$220 tests ; chaque hop inter-modèle \textbf{scellé et re-vérifiable}.
\item \textbf{L'Encadrement} -- sur un run réel de bout en bout, le pré-traitement a \textbf{réduit
l'entrée du gros modèle de $\sim$85\,\%} (132→20 tokens), réponse \textbf{vérifiée SUPPORTED}.
\item \textbf{Essaim médiateur SLM} -- bâti sur le hop \emph{éprouvé} du Relais (les \textbf{deux
tiers déjà là}, testés) ; la chaîne déclarée en est l'extension naturelle.
\item \textbf{La méthode} -- un \textbf{instrument de mesure re-prouvable} et discipliné : résultats
reproductibles, constantes gelées.
\end{itemize}

\end{document}
